\begin{abstract}
Long-context systems commonly retain or reconstruct key--value state at every
layer.  We study a training-free alternative on a frozen backbone: cache one
activation tensor per chunk at depth $j$, retrieve a fixed number of chunks, and
recompute only the upper layers.  \textbf{CoMem} implements this
write--retrieve--read pipeline.  With fixed top-$k$, model read length is
independent of document length and the measured GPU read-time working set remains
approximately flat through 128k tokens.  The external activation store, raw text,
and BM25 index still grow with the document, so this is not a constant-memory
system.

Matched measurements on Qwen3-8B separate the retrieval operating point from the
depth-cache trade-off.  A fixed-$k$ short read provides the working-set benefit.
At the same 6.7k-token read in a 128k memory check, moving from $j{=}0$ to
$j{=}12$ saves about 0.3\,GB, a read-length-scoped quantity rather than a general
memory benefit.  The matched $j{=}12$ read is less accurate than re-forwarding the
same retrieved chunks from $j{=}0$.  In a five-arm decomposition, the contrast
remains significant when selection, decoding, chunk-independent encoding, and
position relocation are held fixed; a residual lower-band context-set difference
prevents interpreting it as a fully deconfounded causal effect.  The relocation
contrast is null at the scored endpoint but not at the raw-generation level.

The preregistered depth sweep has the same direction: pointwise and without
multiplicity correction, every tested $j{\ge}2$ read loses to matched $j{=}0$ on
the same samples, and no deeper point recovers it.  The load-bearing $j{=}12$
comparison remains significant under Bonferroni correction, although the curve is
not strictly monotone.  Because the headline 4k magnitude was selected and depends
on the registered reader convention, the reusable conclusion is the direction,
not a convention-independent effect size.  The result bounds what this frozen,
training-free backbone gives up; it does not bound recovery after training for
cache compressibility or establish superiority over unimplemented prior systems.
\end{abstract}
